\section[Monthly Recursive SVAR Event-Study Atlas (1960--1973)]{Monthly Recursive SVAR Event-Study Atlas\\and Discrete Multipliers (1960--1973)}
\label{app:annual_svar}
\label{app:recursive_svar}

\textbf{B.1} To evaluate structural parameter stability across historical turning points and trace the dynamic transmission mechanism without the confounding distortions of temporal aggregation, Table~\ref{tab:svar_event_comparison} documents the recursive Structural Vector Autoregression (SVAR) impulse response matrix across four expanding estimation windows extended to a 12-month horizon ($h = 0 \dots 12$), accompanied by the Forecast Error Variance Decomposition (FEVD) of monthly inflation.

\textbf{B.2} Following \citet{ChenFangLiu2021}, \citet{Wright2012}, and \citet{KilianLutkepohl2017}, conventional event-study designs relying on static dummy regressions ($y_t = \alpha + \sum \gamma_j D_{j,t} + \varepsilon_t$) suffer from confounding simultaneous policy shocks, anticipatory leakages, and static truncation. In contrast, under the moving average representation of an SVAR ($y_{t+h} - y_{t-1} = \alpha_h + \theta_h u_t + \dots$), projecting identified structural innovations across event-anchored historical windows isolates exogenous external shocks and captures dynamic multi-period propagation. The structural state vector is specified on stationary monthly series:
\begin{equation}
	\mathbf{Y}_t = \big[ \Delta \ln \Theta_t,\; g_{\text{Manuf},t},\; g_{M0,t},\; \pi_t \big]'
	\label{eq:app_b_svar_vector}
\end{equation}
where $\Delta \ln \Theta_t$ is the monthly growth rate of the Real Structural Imbalance ($\Theta_t \equiv M_t / M^{\text{cap}}_t \times 100$, Appendix~\ref{app:archival_codebook}), $g_{\text{Manuf},t}$ is manufacturing output growth, $g_{M0,t}$ is Central Bank base money growth, and $\pi_t$ is monthly consumer inflation. The structural VAR system is formulated as:
\begin{equation}
	\mathbf{A}_0 \mathbf{Y}_t = \mathbf{c} + \sum_{j=1}^p \mathbf{A}_j \mathbf{Y}_{t-j} + \mathbf{u}_t, \quad \mathbf{u}_t \sim \mathcal{N}(0, \mathbf{D})
	\label{eq:app_b_svar_system}
\end{equation}
where $\mathbf{D} = \text{diag}(\sigma_1^2, \sigma_2^2, \sigma_3^2, \sigma_4^2)$ is strictly diagonal, guaranteeing that structural innovations $\mathbf{u}_t$ are mutually orthogonal. Contemporaneous interactions are governed by the structural matrix:
\begin{equation}
	\mathbf{A}_0 = \begin{bmatrix} 1 & 0 & 0 & 0 \\ -a_{21} & 1 & 0 & 0 \\ -a_{31} & -a_{32} & 1 & 0 \\ -a_{41} & -a_{42} & -a_{43} & 1 \end{bmatrix}
	\label{eq:app_b_a0_matrix}
\end{equation}
The zero contemporaneous restrictions ($a_{12} = a_{13} = a_{14} = a_{23} = a_{24} = a_{34} = 0$) are not arbitrary conventions: they are the exact empirical properties verified by the non-parametric Momentary Conditional Independence (MCI) tests of the Peter-Clark Momentary Conditional Independence (PCMCI) algorithm \citep{Runge2019} in Appendix~\ref{app:causal_pcmci} (Table~\ref{tab:pcmci_monthly_gate} and Figure~\ref{fig:dag_pcmci_monthly_gate}). Because the discovered contemporaneous Directed Acyclic Graph (DAG) is strictly lower-triangular and acyclic, the recursive Cholesky factorization $\mathbf{\Sigma} = \mathbf{P} \mathbf{P}'$ (where $\mathbf{P} = \mathbf{A}_0^{-1} \mathbf{D}^{1/2}$) represents the unique data-certified identification of structural innovations. Structural parameters are estimated recursively across four expanding institutional windows: (i) Expansion (up to Dec 1971, $N=143$), (ii) Pre-Lockout (up to Sep 1972, $N=152$), (iii) Bosses' Lockout (up to Aug 1973, $N=163$), and (iv) Coup and Price Deregulation (up to Dec 1973, $N=167$).

\textbf{B.3} The empirical findings in Table~\ref{tab:svar_event_comparison} establish four key results, with statistical significance evaluated via 68\% bootstrap confidence intervals ($B=500$ replications, \citealp{SimsZha1999}): (1) \textbf{Falsification of Monetarism ($H_1$):} in Window 1, base money explains only $0.11\%$ of inflation variance at $h=1$ and $2.81\%$ at $h=12$, while manufacturing ($5.84\%$) and structural imbalance ($4.88\%$) dominate; (2) \textbf{Emergence of Pass-Through (The External Switch):} structural imbalance shocks produce statistically zero price impact in Window 1 ($-0.0549\%$ at $h=0$, 68\% CI $[-0.20\%, +0.07\%]$, zero contained; and $+0.0009\%$ at $h=12$). Under reserve exhaustion in Windows 2--3, the response flips strictly positive and statistically significant: $+0.2337\%$ at $h=0$ (68\% CI $[+0.01\%, +0.43\%]$, zero excluded) in Window 2, and $+0.2685\%$ at $h=0$ (68\% CI $[+0.04\%, +0.48\%]$, zero excluded), persisting at $h=3$ ($+0.1627\%$, 68\% CI $[+0.04\%, +0.25\%]$) and $h=6$ ($+0.0780\%$, 68\% CI $[+0.01\%, +0.13\%]$) in Window 3; (3) \textbf{Defensive Liquidity Surge ($H_3$):} base money expansion accommodates output growth across all windows ($+0.4331\%$, $+0.3221\%$, $+0.2550\%$, $+0.3264\%$ at $h=0$), with variance explanation reaching $8.19\%$ at $h=12$ in Window 3 as the Central Bank refinanced enterprise payrolls under distribution sabotage; and (4) \textbf{6-Month Stabilization Law:} information propagation stabilizes by month 6 across all windows ($\Delta \text{FEVD}_{12-6} < 0.70\%$), while own-price variance accounts for $97.48\%$ in Window 4 under Decree Law 522 decontrol.

\begin{table}[htbp]
\centering
\renewcommand{\arraystretch}{0.85}
\caption{Recursive SVAR Impulse Responses and Variance Decomposition on Observed Empirical Monthly Series (12-Month Horizon)}
\label{tab:svar_event_comparison}
\resizebox{0.95\textwidth}{!}{%
\begin{tabular}{@{}l c rrrr@{}}
\hline\hline
\multicolumn{6}{l}{\textbf{Panel A: Impulse Responses to One-Standard-Deviation Shocks ($h = 0 \dots 12$)}} \\[0.2em]
\hline
Event Window & Horizon & $\pi$ Resp to $\Delta \ln \Theta$ & $g_{M0}$ Resp to $\Delta \ln \Theta$ & $g_{M0}$ Resp to Manuf & $\pi$ Resp to Manuf \\
\hline
Up to 1971 (Expansion) & h=0 & -0.0549\%~{\scriptsize [-0.20, 0.07]} & -0.6471\%*~{\scriptsize [-1.16, -0.19]} & 0.4331\%~{\scriptsize [-0.17, 1.05]} & -0.3555\%*~{\scriptsize [-0.49, -0.20]} \\
 & h=1 & -0.3191\%*~{\scriptsize [-0.43, -0.18]} & 0.4795\%*~{\scriptsize [0.01, 0.96]} & -0.8912\%*~{\scriptsize [-1.30, -0.41]} & -0.2106\%*~{\scriptsize [-0.32, -0.09]} \\
 & h=3 & 0.1032\%*~{\scriptsize [0.04, 0.17]} & -0.6874\%*~{\scriptsize [-1.07, -0.27]} & -0.1934\%~{\scriptsize [-0.50, 0.12]} & 0.0143\%~{\scriptsize [-0.04, 0.08]} \\
 & h=6 & 0.0294\%~{\scriptsize [-0.00, 0.05]} & -0.2188\%*~{\scriptsize [-0.36, -0.07]} & -0.0940\%~{\scriptsize [-0.18, 0.01]} & 0.0187\%*~{\scriptsize [0.00, 0.04]} \\
 & h=12 & 0.0009\%~{\scriptsize [-0.00, 0.00]} & -0.0071\%~{\scriptsize [-0.03, 0.02]} & -0.0165\%~{\scriptsize [-0.03, 0.00]} & 0.0022\%~{\scriptsize [-0.00, 0.00]} \\
\hline
Pre-Lockout (Sept 1972) & h=0 & 0.2337\%*~{\scriptsize [0.01, 0.43]} & -0.6100\%*~{\scriptsize [-1.11, -0.19]} & 0.3221\%~{\scriptsize [-0.20, 0.90]} & -0.4895\%*~{\scriptsize [-0.63, -0.32]} \\
 & h=1 & 0.0843\%~{\scriptsize [-0.13, 0.30]} & 0.5356\%*~{\scriptsize [0.10, 0.95]} & -0.9627\%*~{\scriptsize [-1.33, -0.53]} & -0.2811\%*~{\scriptsize [-0.46, -0.07]} \\
 & h=3 & 0.0923\%~{\scriptsize [-0.00, 0.17]} & -0.4689\%*~{\scriptsize [-0.81, -0.11]} & -0.1776\%~{\scriptsize [-0.45, 0.08]} & -0.1876\%*~{\scriptsize [-0.27, -0.09]} \\
 & h=6 & 0.0384\%*~{\scriptsize [0.00, 0.06]} & -0.1413\%*~{\scriptsize [-0.27, -0.00]} & -0.0965\%~{\scriptsize [-0.20, 0.01]} & -0.0517\%*~{\scriptsize [-0.08, -0.00]} \\
 & h=12 & 0.0038\%~{\scriptsize [-0.00, 0.01]} & -0.0077\%~{\scriptsize [-0.03, 0.01]} & -0.0210\%*~{\scriptsize [-0.04, -0.00]} & -0.0048\%~{\scriptsize [-0.01, 0.00]} \\
\hline
Bosses Lockout (Aug 1973) & h=0 & 0.2685\%*~{\scriptsize [0.04, 0.48]} & -0.3991\%*~{\scriptsize [-0.87, -0.00]} & 0.2550\%~{\scriptsize [-0.17, 0.79]} & -0.3079\%*~{\scriptsize [-0.50, -0.10]} \\
 & h=1 & 0.1117\%~{\scriptsize [-0.13, 0.36]} & 0.6704\%*~{\scriptsize [0.20, 1.08]} & -0.7883\%*~{\scriptsize [-1.14, -0.36]} & -0.0513\%~{\scriptsize [-0.27, 0.17]} \\
 & h=3 & 0.1627\%*~{\scriptsize [0.04, 0.25]} & -0.3776\%*~{\scriptsize [-0.70, -0.05]} & -0.0920\%~{\scriptsize [-0.37, 0.16]} & -0.0861\%~{\scriptsize [-0.21, 0.05]} \\
 & h=6 & 0.0780\%*~{\scriptsize [0.01, 0.13]} & -0.0761\%~{\scriptsize [-0.19, 0.05]} & -0.0586\%~{\scriptsize [-0.16, 0.05]} & -0.0148\%~{\scriptsize [-0.08, 0.06]} \\
 & h=12 & 0.0208\%~{\scriptsize [-0.00, 0.04]} & 0.0115\%~{\scriptsize [-0.02, 0.04]} & -0.0179\%~{\scriptsize [-0.04, 0.01]} & -0.0071\%~{\scriptsize [-0.03, 0.01]} \\
\hline
Coup and Dereg (Dec 1973) & h=0 & 0.2954\%~{\scriptsize [-0.03, 0.57]} & -0.3507\%~{\scriptsize [-0.81, 0.06]} & 0.3264\%~{\scriptsize [-0.15, 0.79]} & 1.0269\%~{\scriptsize [-0.47, 2.52]} \\
 & h=1 & -0.8658\%*~{\scriptsize [-1.29, -0.41]} & 0.6530\%*~{\scriptsize [0.24, 1.05]} & -0.6997\%*~{\scriptsize [-1.11, -0.17]} & -1.1566\%*~{\scriptsize [-1.60, -0.57]} \\
 & h=3 & 0.1605\%~{\scriptsize [-0.10, 0.37]} & -0.6632\%*~{\scriptsize [-0.97, -0.34]} & -0.3198\%*~{\scriptsize [-0.60, -0.02]} & 0.3581\%*~{\scriptsize [0.11, 0.60]} \\
 & h=6 & -0.0147\%~{\scriptsize [-0.11, 0.08]} & -0.1571\%*~{\scriptsize [-0.26, -0.04]} & -0.0937\%~{\scriptsize [-0.20, 0.03]} & 0.0731\%~{\scriptsize [-0.05, 0.21]} \\
 & h=12 & -0.0060\%~{\scriptsize [-0.02, 0.01]} & -0.0072\%~{\scriptsize [-0.02, 0.01]} & -0.0149\%~{\scriptsize [-0.04, 0.00]} & -0.0093\%~{\scriptsize [-0.03, 0.01]} \\
\hline
\\[-0.5em]
\multicolumn{6}{l}{\textbf{Panel B: Forecast Error Variance Decomposition of Monthly Inflation ($\pi_t$) (\%)}} \\[0.2em]
\hline
Event Window & Horizon & Structural Imbalance ($\Delta \ln \Theta$) & Output ($g_{\text{Manuf}}$) & Base Money ($g_{M0}$) & Inflation Own ($\pi_t$) \\
\hline
Up to 1971 (Expansion) & h=1 & 0.12\% & 5.14\% & 0.11\% & 94.64\% \\
 & h=3 & 4.46\% & 5.84\% & 2.47\% & 87.23\% \\
 & h=6 & 4.85\% & 5.83\% & 2.80\% & 86.52\% \\
 & h=12 & 4.88\% & 5.84\% & 2.81\% & 86.46\% \\
\hline
Pre-Lockout (Sept 1972) & h=1 & 0.97\% & 4.28\% & 0.21\% & 94.54\% \\
 & h=3 & 0.70\% & 3.52\% & 2.84\% & 92.94\% \\
 & h=6 & 0.73\% & 3.65\% & 3.93\% & 91.69\% \\
 & h=12 & 0.74\% & 3.65\% & 4.06\% & 91.55\% \\
\hline
Bosses Lockout (Aug 1973) & h=1 & 1.07\% & 1.41\% & 0.63\% & 96.90\% \\
 & h=3 & 0.65\% & 0.76\% & 5.39\% & 93.20\% \\
 & h=6 & 0.75\% & 0.68\% & 7.49\% & 91.07\% \\
 & h=12 & 0.77\% & 0.64\% & 8.19\% & 90.41\% \\
\hline
Coup and Dereg (Dec 1973) & h=1 & 0.18\% & 2.22\% & 0.12\% & 97.48\% \\
 & h=3 & 1.47\% & 4.20\% & 4.49\% & 89.85\% \\
 & h=6 & 1.45\% & 4.28\% & 6.25\% & 88.02\% \\
 & h=12 & 1.44\% & 4.24\% & 6.66\% & 87.67\% \\
\hline
\hline\hline
\multicolumn{6}{p{15.2cm}}{\small Notes: Recursive Structural Vector Autoregression (SVAR) estimated on strictly observed empirical monthly series ($\Delta \ln \Theta_t$, $g_{\text{Manuf},t}$, $g_{M0,t}$, $\pi_t$) using the R \texttt{vars} package. Contemporaneous identification follows the recursive lower-triangular Cholesky ordering authorized by the non-parametric PCMCI DAG (Appendix~\ref{app:causal_pcmci}): external structural imbalance ($\Delta \ln \Theta_t$) is predetermined, manufacturing output ($g_{\text{Manuf},t}$) responds next, Central Bank base money ($g_{M0,t}$) accommodates, and consumer inflation ($\pi_t$) responds contemporaneously to all variables. Panel~A reports responses to a 1 standard-deviation positive structural imbalance shock and a 1 standard-deviation positive manufacturing output shock. Point estimates reported to four decimals; brackets report 68\% bootstrap confidence intervals $[\text{lower}, \text{upper}]$ ($B=500$ replications) computed following \citet{SimsZha1999}. Asterisks ($^*$) indicate parameter responses whose 68\% bootstrap credibility envelope strictly excludes zero, establishing statistical significance. Panel~B reports the 12-month Forecast Error Variance Decomposition (FEVD) of inflation.}
\end{tabular}%
}
\end{table}

